% ID: SEL-D0027-OA-E03
% ESTADO: validada. S028, cotejada con examen y criterios.
\textbf{a)} $P(A^c)=0.4$. Las dos vías disjuntas para que ocurra $B$
son que ocurra junto a $A$ o junto a $A^c$. Por la regla del producto,
$P(A\cap B)=0.6(0.8)=\boxed{0.48}$ y $P(A^c\cap B)=0.4(0.1)=0.04$.
La probabilidad total da $P(B)=0.48+0.04=\boxed{0.52}$.
Por inclusión-exclusión,
$P(A\cup B)=0.6+0.52-0.48=\boxed{0.64}$.

\textbf{b)} Si los sucesos fueran independientes,
$P(B\mid A)$ sería igual a $P(B)$. Como $0.8\ne0.52$,
los sucesos son dependientes.
